\documentclass{article}
\usepackage{graphicx} % Required for inserting images 

\title{bachelor thesis}
\author{Matthias Peter Limmer }
\date{March 2026}

\begin{document}

\maketitle

\section{Introduction}

goal is 30 to 40 sides

\section{Goal}
Interactive Web-Based Visualization of Genomic Language Model Dependencies

Context
A recent study (Silva et al., Nature Genetics 2025) introduced "nucleotide dependency analysis," a method to extract biological grammar such as RNA structures and transcription factor binding sites from Genomic Language Models (gLMs). Currently, these dependencies exist as raw, high-dimensional tensors (Shape N×N×4x4), making them difficult to interpret without specialized scripting.

Project Description
This thesis aims to develop a web-based visualization tool that acts as a microscope for these models. The goal is to create a application where a user can simply upload a raw PyTorch tensor (pickled) along its start coordinates, and the system instantly renders an interactive, navigable heatmap of the genomic region.
The core engineering challenge lies in bridging the gap between raw Machine Learning outputs and fluid web visualization. The student will need to build a pipeline that ingests the uploaded tensor, reduces its dimensionality (aggregating the 4 target nucleotide channels), and converts it into a multi-resolution format (such as .mcool). The frontend will be adapted from the HiGlass framework to support semantic zooming, allowing users to transition from viewing an entire region to inspecting individual nucleotide interactions.

Main Outcomes
Ingestion Pipeline: A backend service (Python) capable of loading pickled tensors (N×N×4x4), handling the aggregation logic (e.g., max-pooling across the 4th dimension), and converting them into multi-resolution tile formats (Cooler/Zarr).
Interactive Viewer: A React-based frontend using HiGlass that renders the processed map. 
It should support:
Deep Zooming: From high-level domain views down to single-nucleotide resolution.
Dynamic Styling: Real-time adjustable color scales (blue-to-red diverging maps) to highlight faint RNA structure signals vs. strong regulatory blocks.
Coordinate Mapping: Visualizing the relative tensor indices (0..N, 0..N) alongside a virtual genomic axis.
Deployment: A containerized application (Docker) ready for lab hosting.
Add an option to display a nucleotide logo track above and left of the heatmap
Given a reconstruction tensor (N,4)
Aggregate the (NxNx4x4) tensor to display the average dependency per row (influence score)
Display the original reference nucleotide sequence
Optionally display the 4x4 odds ratio table indicating base-level dependencies at a pair of positions. This can be done by hovering or clicking on the specific pair. (see Fig. 1C, odds ratio table in  (Silva et al., Nature Genetics 2025))



\section{Tools}
higlass

docker

react (vite)
...



\section{Tool analysis}

\section{Problems}

\section{Solutions}

\section{Manual}

\section{Conclusion}

\section{Sources}
\end{document}
